\dabstract{
Organizations that have accumulated systems over time often depend on information scattered across several sources, which requires manual lookup and consolidation and increases the risk of rework and inconsistencies. This study analyzes InLight, an application developed to unify the information of the energy procurement back office of Light's Regulated Contracting Environment (ACR), and proposes its implementation and evaluation, considering the benefits and impacts of enterprise systems unification. The research, which is applied, exploratory and descriptive, adopts a single case study and combines a theoretical background on information systems, integration, business processes and technology acceptance with a literature review conducted in the Scopus database, which resulted in 11 synthesized studies. Based on this framework, the study diagnoses the current scenario of the area, in which information is spread across the CCEE systems, the Thunders platform, an invoice entry system and control spreadsheets. Describes InLight's modules, architecture and technologies proposes an implementation plan based on the enterprise system life cycle phases and on critical success factors and defines process indicators, a user perception questionnaire and a manager interview guide to evaluate its results. Since the application has not yet been implemented, the expected benefits are presented as working hypotheses. The study concludes that the gains from unification depend on the solution addressing concrete process problems, on the process being redesigned together with the technology, and on users perceiving value in the new way of working.
}
{systems unification; enterprise systems; system implementation; Brazilian electricity sector; regulated electricity market.}